\chapter*{Abstract}
\addcontentsline{toc}{chapter}{Abstract}
This study examines the volatility dynamics of the Indian equity market by estimating and comparing four GARCH-family models — GARCH(1,1), EGARCH(1,1), GJR-GARCH(1,1), and FIGARCH(1,d,1) — on daily Nifty 50 returns over the period January 2010 to December 2024, comprising 3,678 trading day observations.\\
Daily log returns are computed from Nifty 50 closing prices obtained from the National Stock Exchange of India. Preliminary analysis confirms the presence of volatility clustering, fat tails, negative skewness, and significant ARCH effects — all satisfying the preconditions for GARCH estimation. All four models are estimated using Maximum Likelihood under Student-t distributed errors, with the ARMA order of the mean equation selected systematically by the Akaike Information Criterion. A Bai-Perron structural break test applied to squared returns identifies two statistically significant breakpoints at 27 February 2020 and 20 May 2022, delineating three distinct volatility regimes — a pre-COVID stable regime, a COVID crisis regime, and a post-crisis moderate regime. Model performance is evaluated across both in-sample fit using AIC and BIC criteria and out-of-sample forecasting accuracy across four horizons — 1-day, 3-day, 5-day, and 20-day — evaluated over Q1 2025 using rolling window estimation.\\
The in-sample analysis finds that EGARCH(1,1) achieves the best fit by AIC, with a statistically significant leverage effect confirmed across both EGARCH and GJR-GARCH specifications. Volatility persistence is high across all models, with the GARCH(1,1) alpha plus beta sum of 0.9798 approaching unity. The FIGARCH fractional differencing parameter d = 0.484 confirms the presence of long memory in Nifty 50 conditional volatility. The sub-period analysis reveals that volatility persistence increased significantly during the COVID crisis regime and partially normalised in the post-crisis period. A notable divergence is observed between in-sample and out-of-sample performance — FIGARCH(1,d,1), despite ranking last by AIC, consistently achieves the lowest forecasting error across all four horizons, while EGARCH records the highest out-of-sample error. This finding is statistically confirmed by the Diebold-Mariano test. Value at Risk estimates derived from the AIC-optimal model pass the Kupiec Proportion of Failures backtest at both the 1\% and 5\% confidence level.\\
 This study contributes the most comprehensive multi-model, multi-horizon GARCH comparison on Nifty 50 data to date, incorporating structural break analysis across a 15-year sample spanning multiple market crises

